% =====================================================================
%  Portfolioleistung Business Intelligence – Thema 1.2
%  Integration der Quellsysteme eines Dart-Händlers mit Excel Power Query
%  DHBW Ravensburg, WWIBE124, Theoriephase 5 (Okt–Dez 2026)
%
%  Layout nach Handreichung WA v3 Kap. 3.4: Times 12 pt,
%  Ränder 2,5 / 2,0 / 2,5 / 2,5 cm, 1,5-zeilig; Tabellen/Abbildungen/Listings einzeilig.
%  Keine Quellenbelege, kein Literaturverzeichnis, kein KI-Verzeichnis (Vorgabe Kirchberg).
%  Aufbau: Abbildung (Screenshot) -> was wurde gemacht -> Entscheidung + Alternative.
%  Overleaf: Compiler pdfLaTeX, Hauptdatei main.tex. Screenshots nach bilder/ legen.
% =====================================================================
\documentclass[a4paper,12pt]{article}
\usepackage[utf8]{inputenc}
\usepackage[T1]{fontenc}
\usepackage{newtxtext}
\usepackage{newtxmath}
\usepackage[ngerman]{babel}
\usepackage[babel, german=quotes]{csquotes}
\usepackage{graphicx}
\graphicspath{{bilder/}}
\usepackage{geometry}
\geometry{left=2.5cm, right=2.5cm, top=2.5cm, bottom=2.0cm}
\usepackage{setspace}
\usepackage{parskip}
\usepackage{xcolor}
\usepackage{tabularx}
\usepackage{booktabs}
\usepackage{array}
\usepackage{float}
\usepackage{placeins}
\usepackage{enumitem}
\setlist{noitemsep}
\usepackage[font=small, labelfont=bf, labelsep=colon]{caption}
\usepackage{listings}
% Listings-Definition für Power Query M und Grundstil
\definecolor{lstbg}{gray}{0.96}
\definecolor{lstkw}{RGB}{0,70,140}
\definecolor{lstcm}{RGB}{90,110,90}
\definecolor{lststr}{RGB}{140,40,40}

\lstdefinelanguage{PowerQueryM}{
  morekeywords={let, in, each, if, then, else, try, otherwise, type, as, is, and, or, not, meta, section, shared, nullable, null, true, false},
  morekeywords=[2]{Table, Text, List, Record, Date, DateTime, DateTimeZone, Number, Value, Csv, Excel, File, Character, Replacer, JoinKind, QuoteStyle, Splitter, Order, Int64},
  sensitive=true,
  morecomment=[l]{//},
  morecomment=[s]{/*}{*/},
  morestring=[b]",
  alsoletter={.},
}

\lstset{
  basicstyle=\ttfamily\footnotesize,
  keywordstyle=\color{lstkw}\bfseries,
  keywordstyle=[2]\color{lstkw},
  commentstyle=\color{lstcm}\itshape,
  stringstyle=\color{lststr},
  backgroundcolor=\color{lstbg},
  frame=single, framerule=0pt, framesep=4pt,
  breaklines=true, breakatwhitespace=false,
  postbreak=\mbox{\textcolor{gray}{$\hookrightarrow$}\space},
  columns=fullflexible, keepspaces=true,
  showstringspaces=false,
  numbers=left, numberstyle=\tiny\color{gray}, numbersep=6pt,
  tabsize=2,
  captionpos=b,
  literate=
    {ä}{{\"a}}1 {ö}{{\"o}}1 {ü}{{\"u}}1 {Ä}{{\"A}}1 {Ö}{{\"O}}1 {Ü}{{\"U}}1 {ß}{{\ss}}1
    {Ã}{{\~A}}1 {¼}{{\textonequarter}}1 {¶}{{\P}}1 {¤}{{\textcurrency}}1 {Ÿ}{{\"Y}}1 {œ}{{\oe}}1 {Œ}{{\OE}}1
    {→}{{$\rightarrow$}}1 {–}{{--}}1 {„}{{\glqq}}1 {“}{{\grqq}}1 {”}{{\grqq}}1
}
\renewcommand{\lstlistingname}{Listing}
\renewcommand{\lstlistlistingname}{Listingverzeichnis}

\usepackage{tikz}
\usetikzlibrary{positioning, arrows.meta, fit, backgrounds, calc}
\usepackage[hidelinks]{hyperref}
\usepackage[nameinlink, ngerman]{cleveref}

\usepackage{etoolbox}
\setstretch{1.5}
\AtBeginEnvironment{table}{\setstretch{1.0}\selectfont}
\AtBeginEnvironment{figure}{\setstretch{1.0}\selectfont}
\AtBeginEnvironment{lstlisting}{\setstretch{1.0}\selectfont}
\clubpenalty=10000
\widowpenalty=10000

% ---------- Hilfsmakros
\newcommand{\todo}[1]{\textcolor{red!80!black}{\textbf{[TODO: #1]}}}
\newcommand{\pq}[1]{\texttt{#1}}        % Power-Query-Bezeichner
\newcommand{\quelle}[1]{\textsf{#1}}    % Quellsystem-Objekte
\newcommand{\skript}[1]{(Skript Kap.~#1)}
% \abb{datei}{Bildunterschrift}{label}{Breite}: zeigt bilder/datei.png, sonst einen Platzhalter mit dem Dateinamen
\newcommand{\abb}[4]{%
  \begin{figure}[H]\centering
  \IfFileExists{bilder/#1.png}{\includegraphics[width=#4\textwidth]{#1.png}}%
    {\fbox{\parbox{#4\textwidth}{\centering\vspace{1.2cm}\textcolor{gray}{Screenshot \texttt{bilder/#1.png}}\vspace{1.2cm}}}}
  \caption{#2}\label{#3}\end{figure}}
% Anhänge ohne Gliederungsnummer (HR 2.5/2.9)
\newcounter{anhang}
\crefname{anhang}{Anhang}{Anhänge}
\Crefname{anhang}{Anhang}{Anhänge}
\newcommand{\anhang}[2]{\clearpage\refstepcounter{anhang}%
  \section*{Anhang \Alph{anhang}: #1}\label{#2}\addcontentsline{toc}{section}{Anhang \Alph{anhang}: #1}}

\newcommand{\arbeitstitel}{Integration der Quellsysteme eines Dart-Händlers mit Excel Power Query}
\newcommand{\abgabedatum}{27. November 2026}
\hypersetup{pdftitle={\arbeitstitel}, pdfauthor={Steffen Waldvogel}}

% =====================================================================
\begin{document}

% ---------- Deckblatt (ungezählt)
\pagenumbering{alph}
\thispagestyle{empty}
\begin{flushright}\includegraphics[width=3.5cm]{dhbw_logo.png}\end{flushright}
\vspace{2.5cm}
\begin{center}
  {\Large \textbf{\arbeitstitel}} \\[0.6cm]
  {\large Fallbeispiel DartPoint GmbH -- Thema 1.2} \\[1.2cm]
  {\large \textbf{PORTFOLIOLEISTUNG}} \\[0.4cm]
  Schriftliche Ausarbeitung im Modul Business Intelligence \\[1cm]
  des Studienganges Wirtschaftsinformatik: Business Engineering \\
  an der Dualen Hochschule Baden-Württemberg Ravensburg \\[1.5cm]
  von \\[0.8cm]
  Steffen Waldvogel \\[1.5cm]
  Ravensburg, \abgabedatum
\end{center}
\vfill
\begin{tabular}{ll}
Matrikelnummer, Kurs: & 3989184, WWIBE124 \\
Fachsemester:         & 5. Semester \\
Dualer Partner:       & Hokubema Maschinenbau GmbH, Sigmaringen \\
Dozent:               & \todo{Prof. Dr. Kirchberg -- Titel/Vorname prüfen} \\
\end{tabular}

% ---------- Verzeichnisse (römisch)
\newpage
\pagenumbering{roman}
\setcounter{page}{1}
\tableofcontents
\newpage
\listoffigures
\addcontentsline{toc}{section}{Abbildungsverzeichnis}
% Tabellenverzeichnis nur bei >= 3 Tabellen (HR 2.7) – derzeit 3, also aktiv
\listoftables
\addcontentsline{toc}{section}{Tabellenverzeichnis}

% =====================================================================
% TEXTTEIL – Ziel 8–12 Seiten; bei vielen Abbildungen dürfen es mehr sein (Vorgabe Kirchberg)
% =====================================================================
\newpage
\pagenumbering{arabic}
\setcounter{page}{1}

\section{Einleitung}\label{sec:einleitung}

Ein großer Händler für Dart-Zubehör setzt drei Systeme ein: eine ERP-Software für Bestellungen, Zahlungen und Lieferungen, einen Online-Shop und ein CRM. Aufgabe ist es, für diese drei Quellsysteme Tabellen zu entwerfen, die für Auswertungen von Verkäufern mit Bezug zu Kundenmerkmalen relevant sind, und sie mit Excel Power Query in einen gemeinsamen, integrierten Datenbestand zu überführen. Die Quellsysteme sind so zu entwerfen, dass alle im Skript genannten Integrationsprobleme auftreten und gelöst werden.

Als Fallbeispiel dient die fiktive \textbf{DartPoint GmbH}, ein Großhändler mit Kunden in Deutschland, Österreich und der Schweiz. Sie beliefert Dartvereine, Fachhändler und Gastronomie über Außen- und Innendienst und verkauft an Privatkunden über einen eigenen Online-Shop. Die drei Systeme sind historisch gewachsen: ERP seit 2012 (Lexware-ähnlich, on-premise, Windows-1252-Export), Shop seit 2018 (Standardsystem, UTF-8, ISO-Datum), CRM seit 2022 (SaaS, GUIDs, US-Datum, Excel-Export). Die Heterogenität ist damit die realistische Folge von drei Werkzeugen aus drei Jahrzehnten.

Der integrierte Datenbestand soll sechs Fragen beantworten, die jeweils Verkäufer und Kundenmerkmal verknüpfen: Umsatz je Verkäufer nach Kundentyp (F1), nach Land und Region (F2), Saisonverlauf (F3), Produktgruppen (F4), Kanalanteil Direkt vs.\ Shop bei betreuten Kunden (F5) sowie Zahlungs- und Lieferqualität je Verkäufer (F6).

\Cref{sec:quellsysteme} entwirft die Quellsysteme und die Beispieldaten, \cref{sec:zielschema} das Zielschema. \Cref{sec:integration} zeigt die Umsetzung in Power Query Schritt für Schritt anhand von Screenshots: je Schritt, was gemacht wurde, warum so und welche Alternative verworfen wurde. Der Integrationserfolg wird nicht nur plausibilisiert, sondern gegen die bekannte Wahrheit der Datengenerierung gemessen (\cref{sec:erfolgsmessung}). \Cref{sec:zusammenfassung} fasst die Ergebnisse zusammen.

% =====================================================================
\section{Themenbearbeitung}\label{sec:themen}

\subsection{Quellsysteme und Beispieldaten}\label{sec:quellsysteme}

\begin{figure}[H]\centering\small
\resizebox{\textwidth}{!}{%
\begin{tikzpicture}[x=1cm, y=1cm,
  sys/.style={draw, rounded corners=2pt, fill=white, inner sep=5pt, text width=4.9cm, align=left, font=\scriptsize, anchor=north},
  tech/.style={draw=gray!60, fill=gray!8, rounded corners=2pt, inner sep=4pt, text width=4.9cm, align=left, font=\scriptsize, anchor=north},
  ref/.style={<->, gray, dashed}]
  \node[sys] (erp) at (0,0) {%
    \textbf{\sffamily ERP (Lexware-Typ, seit 2012)}\\[2pt]
    \pq{KUNDE} (KdNr, Firma, Name1, Vorname, Strasse, PLZ, Ort, Land, Kundengruppe, VertreterNr, Zahlungsziel, Gesperrt, Bemerkung)\\
    \pq{VERTRETER} (VertreterNr, Name, Gebiet, Austritt)\\
    \pq{AUFTRAG} (AuftragNr, KdNr, Datum, VertreterNr, Status, Waehrung, Lieferart, ShopRef)\\
    \pq{AUFTRAGSPOS} (AuftragNr, Pos, ArtNr, Menge, Einheit, EPreisNetto)\\
    \pq{ZAHLUNG}, \pq{LIEFERUNG}, \pq{ARTIKEL}};
  \node[sys] (shop) at (5.6,0) {%
    \textbf{\sffamily Online-Shop (seit 2018)}\\[2pt]
    \pq{customers} (customer\_id, email, first\_name, last\_name, gender, address, country, is\_club, newsletter, created\_at, erp\_customer\_no, company)\\
    \pq{orders} (order\_id, customer\_id, order\_date, status, total\_gross, shipping\_method\_id)\\
    \pq{order\_items} (order\_id, line\_no, sku, qty, unit\_price\_gross)\\
    \pq{products} (sku, name, category)\\
    \pq{shipping\_methods} (id, name)};
  \node[sys] (crm) at (11.2,0) {%
    \textbf{\sffamily CRM (SaaS, seit 2022)}\\[2pt]
    \pq{Accounts} (AccountId, Firma, Anrede, Titel, Vorname, Nachname, Strasse, Hausnr, PLZ, Ort, Land, Kundentyp, Segment, Betreuer, ERP\_KdNr, Erstkontakt, Newsletter, Status, LetzteAenderung, E-Mail)\\
    \pq{Betreuer} (Name, E-Mail, Rolle, Region)\\
    \pq{Aktivitaet} (ActivityId, AccountId, Typ, Datum, Betreuer)};
  \node[tech] (erpt) at (0,-4.7) {%
    \textbf{.txt}, Windows-1252, Semikolon, CRLF\\ Datum TT.MM.JJJJ, Dezimalkomma\\ Preise netto, Währung je Auftrag\\ Schlüssel: KdNr numerisch, VertreterNr, ArtNr sprechend\\ Vollabzug ohne Änderungsdatum};
  \node[tech] (shopt) at (5.6,-4.7) {%
    \textbf{.csv}, UTF-8, Komma, LF\\ Datum ISO 8601 UTC, Dezimalpunkt\\ Preise brutto, immer EUR (auch CH)\\ Schlüssel: customer\_id alphanumerisch, sku numerisch\\ created\_at/updated\_at vorhanden};
  \node[tech] (crmt) at (11.2,-4.7) {%
    \textbf{.xlsx}, ein Blatt je Tabelle\\ Datum MM/DD/YYYY, teils nur MM/YYYY\\ keine Preise, Freitext vom Vertrieb\\ Schlüssel: AccountId als GUID, Betreuer per Name\\ LetzteAenderung je Account};
  % Querverweise unterhalb
  \draw[ref] (erpt.south) to[out=-90, in=-90, looseness=0.5] node[below, font=\tiny, align=center, pos=0.5] {Shop.erp\_customer\_no (40\,\%)} (shopt.south);
  \draw[ref] (shopt.south) to[out=-90, in=-90, looseness=0.5] node[below, font=\tiny, align=center, pos=0.5] {E-Mail (Shop $\leftrightarrow$ CRM)} (crmt.south);
  \draw[ref] (erpt.south west) ++(0.3,0) to[out=-90, in=-90, looseness=0.35] node[below, font=\tiny, align=center, pos=0.5] {CRM.ERP\_KdNr (60\,\%, mit Tippfehlern)} ($(crmt.south east)+(-0.3,0)$);
\end{tikzpicture}}
\caption{Die drei Quellsysteme der DartPoint GmbH: Tabellen, Formate und die einzigen vorhandenen Querverweise}
\label{fig:quellen}
\end{figure}

Das \textbf{ERP-System} ist das älteste der drei Systeme und verwaltet die kaufmännischen Stamm- und Bewegungsdaten: Kunden, Aufträge, Auftragspositionen, Zahlungen und Lieferungen. Der Export erfolgt als Textdatei im Zeichensatz Windows-1252 mit Semikolon als Trennzeichen, deutschem Datumsformat und Dezimalkomma. Strukturell weist es die Merkmale eines gewachsenen Altsystems auf: Straße und Hausnummer werden in einem gemeinsamen Attribut geführt, der Firmenname ist auf 30 Zeichen begrenzt, die Kundengruppe ist numerisch kodiert (1 bis 4), und fehlende Werte werden durch Ersatzwerte dargestellt (Zahlungsziel $-1$, Sperrdatum 31.12.9999). Informationen, für die kein Attribut vorgesehen ist, werden im Freitextfeld \pq{Bemerkung} erfasst, darunter Sperrvermerke und mehrzeilige Einträge.

Der \textbf{Online-Shop} enthält ausschließlich Daten, die der Besteller selbst erfasst. Die Anschrift liegt in einem einzigen Attribut vor, anstelle eines Kundentyps existiert lediglich das Kennzeichen \pq{is\_club}, und ein Verkäuferbezug fehlt vollständig. Preise werden brutto und durchgängig in Euro geführt, auch für Kunden in der Schweiz; Zeitstempel liegen in koordinierter Weltzeit (UTC) vor. Der Bezug zum ERP-System wird über ein optionales Attribut mit der ERP-Kundennummer hergestellt, das bei 40\,\% der Shop-Kunden gefüllt ist.

Das \textbf{CRM-System} wird vom Vertrieb manuell gepflegt, was sich in der Datenqualität niederschlägt. Vereine sind ohne Vereinskürzel und Rechtsform erfasst, Privatkunden teilweise mit vollständigem Namen im Firmenfeld, und der zuständige Betreuer ist als Name statt als Schlüssel hinterlegt. Der Export erfolgt als Excel-Arbeitsmappe mit US-amerikanischem Datumsformat. Dabei sind Postleitzahlen teilweise als Zahl gespeichert und Umlaute teilweise fehlerhaft kodiert. Die ERP-Kundennummer ist bei 60\,\% der Accounts eingetragen, in Einzelfällen fehlerhaft.

Ein systemübergreifender Schlüssel existiert somit nicht. Die beiden Verweisattribute decken nur einen Teil des Kundenbestands ab, und der Verweis im CRM-System ist selbst fehlerbehaftet. Für die übrigen Kunden stehen zur Identifikation lediglich E-Mail-Adresse, Anschrift und Name zur Verfügung, deren Schreibweise sich zwischen den Systemen unterscheidet.

\begin{table}[H]
  \centering\scriptsize
  \caption{Dieselbe Kundin in den drei Quellsystemen (Kunde 31 der Sollwerte): drei Schlüssel, drei Schreibweisen, drei Kodierungen -- und fünf eingebaute Fehler}
  \label{tab:beispielzeilen}
  \begin{tabular}{@{}l>{\ttfamily}p{3.6cm}>{\ttfamily}p{4.2cm}>{\ttfamily}p{4.2cm}@{}}
    \toprule
    Feld & \normalfont ERP \quelle{KUNDE} & \normalfont Shop \quelle{customers} & \normalfont CRM \quelle{Accounts} \\
    \midrule
    Schlüssel & 10275 & C-00174 & 96ea9404-\dots\ (GUID) \\
    Verweis & -- & erp\_customer\_no = (leer) & ERP\_KdNr = 10275 \\
    Firma & Dar-tLounge \normalfont(Tippfehler) & Dart-Lounge & Dart-Lounge \\
    Name & Dr. Bauer; Claudia \normalfont(Titel im Namen) & Claudia; Bauer; gender=2 & Frau; Claudia; Bauer \\
    Adresse & 18 Bahnhofstr. \normalfont(Hausnr. vorn) & "Bahnhof-Straße 18, 80331 München" & Bahnhofstrasse; 18 \\
    PLZ / Ort & 80331; München & (in address) & 80331; München \\
    Land & D & DE & Deutschland \\
    Kundentyp & Kundengruppe = 3 & is\_club = false & Gastronomie \\
    Newsletter & j \normalfont(statt J/N) & true & X \\
    E-Mail & -- & claudia.bauer98@t-online.de \normalfont(Sammelpostfach) & claudia.bauer98@t-online.de \\
    Telefon & 0751/77093 & 0751.67167 & 0049 751 57931 \\
    Datum & Anlage TT.MM.JJJJ & created\_at ISO/UTC & Erstkontakt 01/20/2021 \\
    \bottomrule
  \end{tabular}
\end{table}

\abb{vn-kunde-6}{Kunde 6 in den drei Quellen: Rohwert je Merkmal und harmonisiertes Ergebnis (Abfrage \pq{BELEG\_Heterogenitaet}). Aus drei Schreibweisen der Goethestraße wird ein Adressschlüssel}{fig:vnkunde}{1.0}
Unter derselben Adresse (Bahnhofstraße 18, München) führen alle drei Systeme außerdem einen zweiten Kunden: den Verein \enquote{Black Arrows München}, der sein Vereinsheim in der Gaststätte hat -- im CRM als \enquote{Black Arrows München} mit Zeichensatzfehler. Die E-Mail-Adresse der Gaststätte verwendet zusätzlich ein weiterer Kunde. Beides sind Fälle, in denen ein an sich starkes Merkmal (Adresse, E-Mail) zwei verschiedene Kunden verbindet.

Die Kundenbestände der drei Systeme überlappen sich, ohne dass ein System die anderen vollständig enthält. Vereine, Händler und Gastronomiebetriebe sind im ERP- und im CRM-System geführt, zu einem großen Teil zusätzlich im Shop. Privatkunden sind nahezu ausschließlich im Shop vorhanden, Interessenten ohne Auftrag ausschließlich im CRM-System, gesperrte Altkunden ausschließlich im ERP-System. Für die Bewegungsdaten gilt Entsprechendes: Ein Teil der Shop-Bestellungen wird in das ERP-System übernommen und liegt damit redundant in beiden Systemen vor. Die Verkäufer sind in ERP- und CRM-System dieselben Personen, werden jedoch über unterschiedliche Schlüssel identifiziert (Vertreternummer beziehungsweise Name).

Zwei Besonderheiten des Dart-Handels erschweren die Objektidentifikation gegenüber einem gewöhnlichen Firmenkundenbestand. Erstens haben Vereine ihr Vereinsheim häufig in einem Gastronomiebetrieb, der selbst Kunde ist, sodass eine Anschrift zwei verschiedenen Kunden zugeordnet ist. Zweitens nutzen Vereine und Gastronomiebetriebe teilweise ein gemeinsames Postfach, sodass auch die E-Mail-Adresse kein eindeutiges Merkmal darstellt.

Die Beispieldaten werden durch ein Python-Skript mit festem Startwert erzeugt (Anhang~B). Zwei Alternativen wurden verworfen: die manuelle Anlage, da sich die Fehlerdichte dabei nicht gezielt steuern lässt, und die unmittelbare Erzeugung durch eine KI, da in diesem Fall nicht bekannt ist, welche Datensätze dasselbe reale Objekt beschreiben. Das Skript erzeugt zunächst 285 eindeutige Kunden und leitet daraus die Datensätze der drei Systeme ab. Anschließend werden die Fehler eingebracht und einzeln protokolliert. Die Fehlerbilder sind dem Skript entnommen und um solche ergänzt, die in manuell gepflegten Kundenstammdaten typischerweise auftreten. Der Umfang ist auf rund 500 Kundendatensätze begrenzt, sodass jedes Fehlerbild mehrfach auftritt und zugleich jeder Einzelfall in Excel nachvollzogen werden kann.

\begin{table}[H]
  \centering\small
  \caption{Verteilung der 285 eindeutigen Kunden auf die Quellsysteme}
  \label{tab:verteilung}
  \begin{tabular}{@{}lrl@{}}
    \toprule
    Präsenz & Anzahl & Kundentypen \\
    \midrule
    ERP + Shop + CRM & 64 & Vereine, Händler, Gastro \\
    ERP + CRM & 28 & Vereine, Händler, Gastro \\
    Shop + CRM & 43 & Vereine, Privat, Gastro \\
    nur CRM (Leads) & 21 & Vereine, Händler, Gastro \\
    nur Shop & 93 & Privat \\
    nur ERP (Altkunden, gesperrt) & 36 & Vereine, Händler, Gastro \\
    \bottomrule
  \end{tabular}
\end{table}
Daraus entstehen 499 Kundensätze: 134 im ERP (inkl.\ 4 Dubletten und 2 Testkunden), 201 im Shop (1 Testkunde) und 164 im CRM (7 Dubletten, 1 Testkunde), dazu 478 ERP-Aufträge und 549 Shop-Bestellungen. Der Generator schreibt zwei Sollwert-Dateien, die nicht Teil der Quellen sind: welcher Kunde unter welcher ID in welchem System steckt, und welcher der 318 eingebauten Fehler in welchem Datensatz sitzt. Von den 135 überlappenden Kunden haben 40\,\% eine identische Adresse, 40\,\% eine abweichende Schreibweise und 20\,\% sind umgezogen, wobei das CRM die jüngere Adresse führt.

% =====================================================================
\subsection{Integrierter Datenbestand}\label{sec:zielschema}

Ziel der Integration ist ein gemeinsamer Datenbestand, auf dem sich die sechs Fragestellungen ohne Rückgriff auf die Quellsysteme beantworten lassen. Er besteht aus einer Tabelle der Verkäufe auf Ebene der Auftragsposition sowie je einer Tabelle für Kunden, Verkäufer und Produkte (\cref{fig:zielschema}). Da die Quellsysteme keinen gemeinsamen Kundenschlüssel besitzen, erhält jeder Kunde einen künstlichen Schlüssel; eine Mapping-Tabelle hält fest, welcher Quellschlüssel auf welchen Zielschlüssel abgebildet wurde. Abweichende Anschriften eines Kunden werden nicht überschrieben, sondern mit Gültigkeitszeitraum geführt (\cref{tab:adresse}), sodass Umsätze der zum Auftragsdatum gültigen Region zugeordnet werden.

\begin{figure}[H]\centering\small
\resizebox{0.95\textwidth}{!}{%
\begin{tikzpicture}[x=1cm, y=1cm, >=Latex,
  ent/.style={draw, rounded corners=2pt, fill=white, inner sep=4pt, text width=3.6cm, align=left, font=\scriptsize},
  fact/.style={ent, fill=gray!12, text width=4.2cm},
  key/.style={font=\scriptsize\bfseries}]
  \node[fact] (f) at (0,0) {%
    \textbf{\sffamily FAKT\_VERKAUF} (1 Zeile je Position)\\[2pt]
    \textbf{KundeKey, VerkaeuferKey, ProduktKey}\\
    AuftragID, Pos, Quelle (ERP/Shop), Kanal (Direkt/Online)\\
    Auftragsdatum, Jahr, Quartal, Monat, Saison, Saisonjahr\\
    LandISO, Region, PLZ \textit{(zum Auftragsdatum)}\\
    Kundentyp, Stamm\_VertreterNr, Produktgruppe\\
    Status, Lieferart, MengeStk, NettoEUR\\
    Lieferdauer\_Tage, Zahlungsverzug\_Tage};
  \node[ent, anchor=south] (k) at (-5.6,0.4) {%
    \textbf{\sffamily DIM\_KUNDE}\\[2pt]
    \textbf{KundeKey} (Surrogat)\\
    Firma, Anrede, Titel, Vorname, Nachname\\
    Kundentyp (Verein/Händler/Gastro/Privat)\\
    LandISO, Newsletter, Erstkontakt, Kontaktstatus (Kunde/Lead/gesperrt)\\
    VertreterNr, Betreuer\\
    ERP\_KdNr, ShopID, AccountId, Quellen, AnzahlQuellen};
  \node[ent, anchor=north] (a) at (-5.6,-0.4) {%
    \textbf{\sffamily DIM\_KUNDE\_ADRESSE} (historisiert)\\[2pt]
    \textbf{KundeKey, Version}\\
    Strasse, Hausnr, PLZ, Ort, LandISO, Region\\
    GueltigVon, GueltigBis, CurrentFlag};
  \node[ent, anchor=south] (v) at (5.6,0.4) {%
    \textbf{\sffamily DIM\_VERKAEUFER}\\[2pt]
    \textbf{VerkaeuferKey} (= VertreterNr, \enquote{ONLINE} als Pseudo-Verkäufer)\\
    Nachname, Vorname, Rolle (Außen-/Innendienst/Online)\\
    Gebiet, Region, E-Mail, Austritt, Aktiv};
  \node[ent, anchor=north] (p) at (5.6,-0.4) {%
    \textbf{\sffamily DIM\_PRODUKT}\\[2pt]
    \textbf{ProduktKey}\\
    ArtNr (ERP), Shop\_sku, Bezeichnung, Shop\_name\\
    Produktgruppe, EKPreis\\
    Flag\_Gruppe\_abweichend, Flag\_ohne\_Shop};
  \node[ent, fill=gray!8, above=7mm of k] (m) {%
    \textbf{\sffamily MAP\_SCHLUESSEL}\\[2pt]
    Quelle, \textbf{QuellID} $\rightarrow$ \textbf{KundeKey}\\
    Matchmethode (ERP-Schlüssel / E-Mail / Adresse / Fuzzy / Dublette)};
  \draw[->] (k.east) -- node[sloped, above, font=\tiny] {1 : n} (f.north west);
  \draw[->] (v.west) -- node[sloped, above, font=\tiny] {1 : n} (f.north east);
  \draw[->] (p.west) -- node[sloped, above, font=\tiny] {1 : n} (f.south east);
  \draw[->] (k) -- node[right, font=\tiny] {1 : n (Versionen)} (a);
  \draw[->, dashed] (a.east) -- node[sloped, below, font=\tiny, fill=white, inner sep=1pt] {Adresse zum Auftragsdatum} (f.south west);
  \draw[->, dashed] (m) -- node[right, font=\tiny] {n : 1} (k);
\end{tikzpicture}}
\caption{Integrierter Datenbestand: Verkäufe, Kunden, Verkäufer und Produkte mit Schlüssel-Mapping und Adressversionen (fett: Schlüssel)}
\label{fig:zielschema}
\end{figure}

\begin{table}[H]
  \centering\small
  \caption{\pq{DIM\_KUNDE\_ADRESSE}: Kunde 15 (DC Black Arrows Bregenz) mit zwei Adressversionen -- Gültigkeitszeitraum statt Entscheidung, welche Quelle recht hat}
  \label{tab:adresse}
  \begin{tabular}{@{}rllllllc@{}}
    \toprule
    KundeKey & Version & Straße & PLZ & Ort & GueltigVon & GueltigBis & CurrentFlag \\
    \midrule
    15 & 1 & Schillerstr.\ 89 & 6900 & Bregenz & 01.01.1900 & 12.01.2024 & FALSCH \\
    15 & 2 & Am Sportplatz 35 & 4020 & Linz & 13.01.2024 & -- & WAHR \\
    \bottomrule
  \end{tabular}
\end{table}


% =====================================================================
\subsection{Import und Harmonisierung}\label{sec:integration}\label{sec:aufbau}\label{sec:syntax}\label{sec:schema}

\begin{figure}[H]\centering\small
\resizebox{\textwidth}{!}{%
\begin{tikzpicture}[>=Latex, x=1cm, y=1cm,
  box/.style={draw, rounded corners=1pt, fill=white, minimum height=7mm, inner sep=3pt, font=\ttfamily\scriptsize, align=center},
  layer/.style={draw=gray!60, fill=gray!8, rounded corners=2pt, inner sep=3pt},
  lab/.style={font=\sffamily\scriptsize\bfseries, text=gray!70!black, anchor=east, align=right}]
  % Zeile 1: Rohdaten und Hilfen
  \node[box] (k) at (0,0) {KUNDE, AUFTRAG, \dots};
  \node[box] (c) at (3.6,0) {customers, orders, \dots};
  \node[box] (a) at (7.2,0) {Accounts, Betreuer};
  \node[box, fill=gray!15] (h) at (10.6,0) {fnClean, fnNorm*, fnStrasse, fnTel,\\mapLand, mapKundentyp, Parameter};
  % Zeile 2: Harmonisierung
  \node[box] (kh) at (0,-1.4) {KUNDE\_h};
  \node[box] (sh) at (3.6,-1.4) {SHOP\_h};
  \node[box] (ch) at (7.2,-1.4) {CRM\_h};
  % Zeile 3: Kundenintegration
  \node[box] (kk) at (0,-2.8) {KUNDE\_k,\\ERP\_Dubletten};
  \node[box] (sk) at (3.6,-2.8) {SHOP\_k};
  \node[box] (ck) at (7.2,-2.8) {CRM\_k};
  \node[box] (map) at (10.6,-2.8) {MAP\_SCHLUESSEL};
  % Zeile 4: Zielschema
  \node[box] (dk) at (0,-4.3) {DIM\_KUNDE,\\DIM\_KUNDE\_ADRESSE};
  \node[box] (dp) at (3.6,-4.3) {DIM\_PRODUKT,\\DIM\_VERKAEUFER};
  \node[box] (fv) at (7.2,-4.3) {FAKT\_ERP + FAKT\_SHOP\\$\rightarrow$ FAKT\_VERKAUF};
  % Zeile 5: Prüfung
  \node[box, fill=gray!15] (cs) at (0,-5.7) {CHECK\_Sollwerte,\\CHECK\_Pruefliste};
  \node[box, fill=gray!15] (cf) at (7.2,-5.7) {CHECK\_Fakten};
  \node[box, fill=gray!15] (pv) at (10.6,-5.7) {Pivots P1--P6};
  % Kanten
  \draw[->] (k) -- (kh); \draw[->] (c) -- (sh); \draw[->] (a) -- (ch);
  \draw[->] (kh) -- (kk); \draw[->] (sh) -- (sk); \draw[->] (ch) -- (ck);
  \draw[->] (kk) -- node[above, font=\tiny] {Basis} (sk);
  \draw[->] (kk.south east) to[out=-20, in=200] (ck.south west);
  \draw[->] (ck) -- (sk);
  \draw[->] (sk) -- (dp); \draw[->] (kk) -- (dk); \draw[->] (ck.east) -- (map.west);
  \draw[->] (map) -- (fv); \draw[->] (dk) -- (fv); \draw[->] (dp) -- (fv);
  \draw[->] (dk) -- (cs); \draw[->] (fv) -- (cf); \draw[->] (fv) -- (pv);
  \begin{scope}[on background layer]
    \node[layer, fit=(k)(c)(a)(h)] (L1) {};
    \node[layer, fit=(kh)(sh)(ch)] (L2) {};
    \node[layer, fit=(kk)(sk)(ck)(map)] (L3) {};
    \node[layer, fit=(dk)(dp)(fv)] (L4) {};
    \node[layer, fit=(cs)(cf)(pv)] (L5) {};
  \end{scope}
  \node[lab] at (-1.75,0) {Rohdaten,\\Hilfen};
  \node[lab] at (-1.75,-1.4) {Harmoni-\\sierung};
  \node[lab] at (-1.75,-2.8) {Kunden-\\integration};
  \node[lab] at (-1.75,-4.3) {Ziel-\\schema};
  \node[lab] at (-1.75,-5.7) {Prüfung};
\end{tikzpicture}}
\caption{Schichten der Abfragen in der Mappe (vereinfachte Abhängigkeiten; die Editor-Ansicht \enquote{Abfrageabhängigkeiten} ist mit 60 Knoten nicht lesbar)}
\label{fig:abhaengigkeiten}
\end{figure}

Die Abfragen sind in Schichten aufgebaut (\cref{fig:abhaengigkeiten}). Die Rohdaten-Abfragen enthalten ausschließlich den Import; jede weitere Schicht verweist auf die vorhergehende, ohne sie zu verändern, sodass sich jeder Zwischenstand einzeln prüfen lässt. Zuordnungen wie Länder, Kundentypen, Regionen und Mehrwertsteuersätze liegen als eigene Mapping-Tabellen vor, veränderliche Größen wie Quellordner, Wechselkurs und Ähnlichkeitsschwelle als Parameter. Die Alternative, diese Werte in den Abfragecode zu schreiben, wurde verworfen, da Änderungen dann an mehreren Stellen nachzuziehen wären. Die Quellen werden nacheinander integriert, beginnend mit dem ERP-System, da es die Verkäufer und die auftragsbezogenen Stammdaten führt; es folgen CRM-System und Shop. Eine Integration aller drei Quellen in einem Schritt wurde verworfen, da sich Zuordnungsfehler dann keiner Quelle mehr zuordnen lassen.

\abb{pq-import-erp}{Import der ERP-Datei: Codepage 1252, Semikolon, Gebietsschema de-DE -- und das Problem eingebetteter Zeilenumbrüche}{fig:import}{0.9}

\abb{vn-1-format}{Vorher und nachher, Gruppe 1: Zeichen und Format. Je Fehlerbild ein Datensatz aus dem Fehlerprotokoll des Generators (Abfrage \pq{BELEG\_Fehlerbilder})}{fig:vn1}{1.0}
% Screenshot: Csv.Document-Dialog bzw. Vorschau mit zerrissenen Zeilen (QuoteStyle.None) vs. korrekt (QuoteStyle.Csv): 124 Zeilen -> 90 Datensätze

Die drei Exporte unterscheiden sich in Zeichensatz, Trennzeichen, Dezimaltrennzeichen und Datumsformat (\cref{fig:import}). Jede Typumwandlung erfolgt deshalb mit ausdrücklich angegebenem Gebietsschema. Die automatische Typerkennung von Power Query wurde verworfen, da sie das Gebietsschema des ausführenden Rechners verwendet und ein Datum wie 03/04/2023 je nach Einstellung unterschiedlich deutet. Zeitstempel des Shops werden vor der Datumsbildung von UTC in mitteleuropäische Zeit umgerechnet, da eine Bestellung kurz vor Mitternacht sonst dem falschen Tag zugeordnet würde; die Sommerzeit bleibt dabei unberücksichtigt. Mehrzeilige Einträge im Bemerkungsfeld des ERP-Systems erfordern die Auswertung der Anführungszeichen beim Import, da ein Datensatz andernfalls in mehrere Zeilen zerfällt.

\begin{lstlisting}[language=PowerQueryM, caption={Hilfsfunktion \pq{fnClean}: Zeichensatzfehler, unsichtbare Zeichen, Mehrfach-Leerzeichen und Platzhalter}, label={lst:fnclean}]
(s as nullable text) as nullable text =>
let
    moji = {{"ü","ü"}, {"ö","ö"}, {"ä","ä"}, {"ß","ß"}, {"Ãœ","Ü"}, {"Ö","Ö"}, {"Ä","Ä"}},
    m0 = if s = null then null else List.Accumulate(moji, s, (acc, p) => Text.Replace(acc, p{0}, p{1})),
    t0 = if m0 = null then null else Text.Clean(m0),
    t1 = if t0 = null then null else Text.Replace(t0, Character.FromNumber(160), " "),
    t2 = if t1 = null then null else Text.Remove(t1, {Character.FromNumber(173), Character.FromNumber(8203)}),
    t3 = if t2 = null then null else Text.Combine(List.Select(Text.Split(t2, " "), each _ <> ""), " "),
    t4 = if t3 = null or t3 = "" or List.Contains({"-", "?", "k.a.", "xxx", "n/a", "unbekannt"}, Text.Lower(t3)) then null else t3
in
    t4
\end{lstlisting}

Vor allen feldbezogenen Regeln wird jede Textspalte durch die Funktion \pq{fnClean} bereinigt (\cref{lst:fnclean}). Sie ersetzt fehlerhaft kodierte Umlaute, entfernt Steuerzeichen und nicht sichtbare Zeichen, reduziert Mehrfach-Leerzeichen und setzt Platzhalter wie \enquote{-} oder \enquote{k.A.} auf einen leeren Wert (\cref{fig:vn1}). Eine zentrale Funktion wurde der Bereinigung je Spalte vorgezogen, da allein die CRM-Tabelle 19 Textspalten enthält und nicht sichtbare Zeichen jeden späteren Schlüsselvergleich scheitern lassen. Die Ersetzungsliste deckt nur deutsche Umlaute ab; andere fehlerhaft kodierte Zeichen bleiben unverändert.

Gleiche Sachverhalte sind in den Quellen unterschiedlich kodiert: die Newsletter-Einwilligung als J/N, true/false und X/leer, das Land als Kürzel, ISO-Code und ausgeschriebener Name, der Kundentyp als Ziffer, als Kennzeichen und als Text. Die Vereinheitlichung erfolgt über Mapping-Tabellen (\cref{fig:mapping}) und für Ja/Nein-Felder über die Funktion \pq{fnJN}; Ersatzwerte für fehlende Angaben ($-1$, 31.12.9999) werden zu leeren Werten. Bedingte Spalten wurden als Alternative verworfen, da jede neue Variante eine Änderung des Abfragecodes erfordert, während eine Mapping-Tabelle lediglich um eine Zeile ergänzt wird. Ein nicht zuordenbarer Wert bleibt leer (\cref{fig:vn2}). Beim Kundentyp ist nur eine Abbildung auf den kleinsten gemeinsamen Nenner möglich: Der Shop unterscheidet lediglich Vereine von sonstigen Kunden, und die im CRM-System getrennten Fach- und Großhändler werden zu \enquote{Händler} zusammengefasst. Gleich benannte Attribute mit unterschiedlicher Bedeutung, etwa \pq{Status} als Auftragsstatus im ERP-System, als Bestellstatus im Shop und als Kontaktstatus im CRM-System, werden bei der Übernahme eindeutig benannt.


\abb{pq-mapping}{Mapping-Tabelle \pq{mapLand}: die Sollkodierung je Quelle und alle angetroffenen Varianten}{fig:mapping}{0.9}

\abb{vn-2-kodierung}{Vorher und nachher, Gruppe 2: Kodierungen. Fehlt das Land und ist die PLZ vierstellig, bleibt der Fall offen und steht auf der Prüfliste}{fig:vn2}{1.0}
\abb{pq-strasse}{\pq{KUNDE\_h}: PLZ aus dem Ortsfeld und aus dem Länderpräfix, Straße und Hausnummer getrennt, vertauschte Reihenfolge erkannt}{fig:strasse}{0.95}

\abb{vn-3-adresse}{Vorher und nachher, Gruppe 3: Adresse. Was sich nicht sicher reparieren lässt, erhält keinen Adressschlüssel}{fig:vn3}{1.0}
% Screenshot: KUNDE_h im Editor, Spalten PLZ_norm, Ort, LandISO, Strasse, Hausnr, Flag_Postfach, Flag_Vertauscht, Befunde

Die Fehlerbilder der Beispieldaten gehen über reine Formatfehler hinaus. Ein in manuell gepflegten Stammdaten verbreiteter Fehlertyp ist der \emph{Feldmissbrauch}: der richtige Wert im falschen Attribut oder Freitext an einer Stelle, an der ein kodierter Wert erwartet wird. \Cref{tab:fehlerbilder} zeigt die eingebauten Fehlerbilder mit Fallzahl und Behandlung.

\begin{table}[H]
  \centering\scriptsize
  \caption{Eingebaute Fehlerbilder in den Kundenstammdaten (Auswahl), Fallzahl laut Fehlerlog und Behandlung in Power Query}
  \label{tab:fehlerbilder}
  \begin{tabular}{@{}>{\raggedright\arraybackslash}p{4.4cm}r>{\raggedright\arraybackslash}p{7.3cm}c@{}}
    \toprule
    Fehlerbild (Beispiel) & Fälle & Behandlung & Klasse \\
    \midrule
    E-Mail mit Transportmüll (\texttt{mailto:}, \texttt{Name <\dots>}, \texttt{(at)}) & 42 & \pq{fnNormEmail} & A \\
    Telefon mit Text oder Durchwahl (\texttt{0751/5512 -16 Hr.\ Vogel}) & 27 & Ziffern bis zum ersten Buchstaben $\rightarrow$ E.164; braucht gültiges Land, Rest auf die Prüfliste & H \\
    Land als Variante (\texttt{BRD}, \texttt{Germany}, \texttt{AUT}) & 20 & Mapping-Tabelle \pq{mapLand}; unbekannter Wert bleibt null & A \\
    PLZ als Zahl aus Excel, führende Null fehlt & 16 & Text, in Deutschland auf fünf Stellen auffüllen & A \\
    Zeichensatzfehler (\texttt{München}) & 12 & Ersetzungsliste in \pq{fnClean} & A \\
    PLZ mit Länderpräfix (\texttt{D-88212}) & 10 & Präfix abtrennen; füllt ein fehlendes Land & A \\
    Name nach 30 Zeichen abgeschnitten & 7 & Namensprüfung über gemeinsamen Anfang statt Gleichheit & H \\
    Hausnummer vor Straße, Felder vertauscht & 7 & Rücktausch-Regel in ERP und CRM & A \\
    Betreuer nicht auflösbar (\texttt{S.\ Huber}) & 7 & Abgleich über normalisierten Nachnamen, Rest Prüfliste & H \\
    Privatkunde nur im Firmenfeld (\texttt{Müller Hans}) & 6 & erstes Wort zählt für den Abgleich; Feldkorrektur manuell & H \\
    Tippfehler im Namen (\texttt{Dar-tLounge}) & 6 & Fuzzy-Kandidat, bestätigt durch Adresse & H \\
    PLZ im Ortsfeld (\texttt{88212 Ravensburg}) & 5 & führende Ziffern übernehmen, wenn PLZ leer & A \\
    Zwei E-Mails in einem Feld, Leerzeichen, Bemerkung im Feld & 12 & teilen, bereinigen, ohne \texttt{@} $\rightarrow$ null & A \\
    Sperrvermerk im Freitext (\texttt{INSOLVENT}) & 4 & Schlüsselwortliste setzt Kontaktstatus \enquote{gesperrt} & A \\
    Postfach im Straßenfeld, Hausnummer-Bereich (\texttt{2-5}) & 6 & erkennen; Postfach erhält keinen Adressschlüssel & H \\
    PLZ fehlt, ungültig oder passt nicht zum Land & 7 & nur erkennen und melden & M \\
    E-Mail ohne \texttt{@} & 2 & nur erkennen und melden & M \\
    Ansprechpartner im Firmenfeld (\texttt{Dart-Shop / Hr.\ Koch}) & 2 & nur erkennen und melden & M \\
    \bottomrule
    \multicolumn{4}{@{}l}{A = automatisch behoben, H = hybrid (Regel mit zweiter Bedingung oder Vorschlag), M = manuell (erkennen und melden)}
  \end{tabular}
\end{table}

\abb{vn-4-kontakt}{Vorher und nachher, Gruppe 4: Kontaktdaten}{fig:vn4}{1.0}

\abb{vn-5-feldmissbrauch}{Vorher und nachher, Gruppe 5: Feldmissbrauch. Das Erkennen gelingt automatisch, die Korrektur des Feldes muss manuell überprüft und entschieden werden}{fig:vn5}{1.0}

Die feldbezogenen Regeln folgen einer festen Reihenfolge: Postleitzahl und Ort, Land, Straße und Hausnummer, Telefon. Die Reihenfolge ist sachlich bedingt, da spätere Regeln die Ergebnisse früherer voraussetzen: Das Länderpräfix einer Postleitzahl liefert ein fehlendes Land, und die Normierung der Telefonnummer benötigt ein gültiges Land für die Vorwahl (\cref{fig:strasse,fig:vn3,fig:vn4}). Fehlt das Land, lässt eine fünfstellige Postleitzahl im Vertriebsgebiet nur Deutschland zu; bei einer vierstelligen bleibt das Land offen, da Österreich und die Schweiz in Frage kommen. Die Trennung von Straße und Hausnummer erfolgt an der ersten Ziffer und erkennt vorangestellte Hausnummern, Hausnummernbereiche und Postfächer. Straßennamen, die selbst eine Zahl enthalten, würden dabei falsch getrennt; sie kommen in den Beispieldaten nicht vor.

Für alle Regeln gilt der Grundsatz, keinen Wert zu schätzen. Ein Wert, der sich nicht sicher ableiten lässt, bleibt leer, und der Datensatz erhält einen Befund für die Prüfliste (\cref{sec:pruefliste}). Bei Feldmissbrauch wird der Fehler erkannt, die Korrektur des Feldes muss jedoch manuell überprüft und entschieden werden (\cref{fig:vn5}). Die Alternative, mängelbehaftete Datensätze zu verwerfen, wurde nicht gewählt, da damit auch die zugehörigen Umsätze entfielen. Verworfen werden ausschließlich Testkunden.

\subsection{Kundenintegration und Erfolgsmessung}\label{sec:linkage}\label{sec:erfolgsmessung}

\todo{Hier begründen: künstlicher Schlüssel statt ERP-Kundennummer (157 von 285 Kunden stehen nicht im ERP); Anschrift als Version mit Gültigkeit statt Priorisierung einer Quelle; übrige Stammattribute: erster vorhandener Wert in der Reihenfolge ERP, CRM, Shop.}

\abb{pq-kaskade}{Match-Kaskade in \pq{CRM\_k}: ERP-Schlüssel mit Bestätigung, Adresse mit Namensbestätigung, Fuzzy mit fachlicher Nachregel, neu}{fig:kaskade}{0.95}
% Screenshot: Spalten Key_Schluessel, Key_Adr, Key_Fuzzy, Matchmethode in CRM_k; alternativ Gruppierung nach Matchmethode

Die drei Quellen haben keinen gemeinsamen Schlüssel: ERP und CRM teilen \pq{ERP\_KdNr} nur bei 60\,\% der Accounts und teils fehlerhaft, Shop und CRM teilen die E-Mail, das ERP führt keine E-Mail. Je Quelle wird deshalb eine Kaskade durchlaufen, die den sichersten verfügbaren Abgleich zuerst versucht und die Matchmethode je Zeile protokolliert (\cref{tab:kaskade}).

\begin{table}[H]
  \centering\small
  \caption{Match-Kaskade je Quelle und Trefferzahlen (Messung 8)}
  \label{tab:kaskade}
  \begin{tabular}{@{}c>{\raggedright\arraybackslash}p{6.2cm}>{\raggedright\arraybackslash}p{3.3cm}>{\raggedright\arraybackslash}p{3.5cm}@{}}
    \toprule
    Stufe & Kriterium & CRM $\rightarrow$ ERP & Shop $\rightarrow$ ERP/CRM \\
    \midrule
    0 & Vorfilter: Testkunden entfernen, Dubletten je Quelle zusammenlegen & 1 Test, 7 Dubletten & 1 Test \\
    1 & ERP-Schlüssel, bestätigt durch Namen oder Adresse & 54 & 21 \\
    2 & E-Mail; bei mehreren Kunden je Adresse entscheidet der Name & entfällt (ERP ohne E-Mail) & 83 \\
    3 & Adressschlüssel und Namensbestätigung & 23 & 3 \\
    4 & Fuzzy auf normalisierten Namen, Nachregel entscheidet & 8 & 0 \\
    -- & neu (eigener \pq{KundeKey}) & 71 & 93 \\
    \bottomrule
  \end{tabular}
\end{table}

\abb{vn-6-schluessel}{Vorher und nachher, Gruppe 6: Datensätze und Schlüssel. Testkunden werden verworfen, Dubletten und Schlüsselfehler löst die Kaskade}{fig:vn6}{1.0}

\todo{Entscheidungen, je mit Alternative: (1) Reihenfolge ERP $\rightarrow$ CRM $\rightarrow$ Shop. (2) Adressschlüssel = Land + PLZ + normierte Straße + erste Hausnummer, ohne Namen. (3) Kein Merkmal zählt allein: Fremdschlüssel braucht Namen oder Adresse (3 Tippfehler im Schlüssel, 7 abgeschnittene Namen -- deshalb gemeinsamer Anfang statt Gleichheit), E-Mail braucht bei Mehrfachverwendung den Namen (5 Sammelpostfächer), Adresse braucht den Namen (8 Vereinsheime). (4) Belegt-Regel für den Fall Dart-Lounge Bregenz/Wien. (5) Alternative Punktemodell (additive Gewichte, Schwellen für Hinweis und Warnung, manuelle Bestätigung) -- verworfen, weil der Fuzzy-Join keinen Ähnlichkeitswert liefert und Gewichte ohne Trainingsdaten geschätzt wären; die Kaskade ist je Stufe messbar.}

\todo{Beobachtung: Fuzzy-Join vervielfacht Zeilen (122 statt 115) und liefert keinen Ähnlichkeitswert. Nachregel: Verein $\rightarrow$ nur identischer normierter Name (Ort steckt im Namen); Händler/Gastro $\rightarrow$ gleiche PLZ oder identischer Name bei gleichem Ort. Alternative \enquote{Schwellwert hochdrehen} -- verworfen, siehe Messung.}

\abb{pq-check}{\pq{CHECK\_Sollwerte}: vergebene Keys gegen die bekannte Wahrheit -- zerrissene und verschmolzene Kunden}{fig:check}{0.55}

Die Integration wird gemessen statt nur plausibilisiert: \pq{CHECK\_Sollwerte} vergleicht die vergebenen \pq{KundeKey}s mit der Sollwerttabelle der Datengenerierung und zählt \emph{zerrissene} Kunden (ein echter Kunde, mehrere Keys) und \emph{verschmolzene} Keys (ein Key, mehrere echte Kunden). \Cref{tab:messreihe} zeigt die Messreihe.

\begin{table}[H]
  \centering\small
  \caption{Messreihe der Kundenintegration auf dem Datenstand mit 285 Kunden und 499 Kundensätzen (Excel, \pq{CHECK\_Sollwerte})}
  \label{tab:messreihe}
  \begin{tabular}{@{}c>{\raggedright\arraybackslash}p{9.4cm}rr@{}}
    \toprule
    Nr. & Stand der Kaskade & zerrissen & verschmolzen \\
    \midrule
    6 & jedes Merkmal zählt allein: Schlüssel bei gleichem Namen, E-Mail, Adresse, Fuzzy mit Nachregel & 17 & 9 \\
    7 & abgesichert: Schlüssel, E-Mail und Adresse brauchen eine Bestätigung & 7 & 1 \\
    8 & wie 7, zusätzlich Belegt-Regel & \textbf{7} & \textbf{0} \\
    \bottomrule
  \end{tabular}
\end{table}
Die Messungen 1 bis 5 liefen auf einem ersten Datenstand mit 200 Kunden ohne Feldmissbrauch. Sie führten von 12 Verschmelzungen (Fuzzy ohne fachliche Regel) auf 0 und zeigten, dass der Fuzzy-Schwellwert zwischen 0,7 und 0,9 wirkungslos wird, sobald eine fachliche Nachregel über die Kandidaten entscheidet.

\todo{Drei Absätze. (1) Messung 6: dieselbe Kaskade, die auf dem ersten Datenstand 0 Verschmelzungen hatte, verschmilzt 9 Keys, sobald Sammelpostfächer und Vereinsheime vorkommen -- eine fehlerfreie Messung gilt nur für die Fehlerbilder, die im Datensatz stecken. (2) Messung 7 und 8: Bestätigung statt Einzelmerkmal; der verbleibende Fall (Dart-Lounge zieht von Bregenz nach Wien, wo es eine zweite gibt) und die Belegt-Regel. (3) Precision vor Recall: die 7 Zerrissenen sind Umzieher ohne bestätigenden Schlüssel, bewusst getrennt, auf der Prüfliste. Methodik: Regeln vorab in einem Python-Nachbau geprüft, Abweichung in Messung 6 (16/10 statt 17/9) durch unterschiedliche Ähnlichkeitsmaße; maßgeblich sind die Excel-Werte. Laufzeit: 8 Minuten je Messung, 20 Minuten Gesamtlauf.}

\subsection{Prüfliste: automatisch, hybrid, manuell}\label{sec:pruefliste}

\abb{pq-pruefliste}{\pq{CHECK\_Pruefliste}: Restbefunde nach der Bereinigung, eine Zeile je Datensatz und Befund}{fig:pruefliste}{0.9}

Was die Regeln nicht sicher beheben können, wird nicht geraten, sondern gemeldet. \pq{CHECK\_Pruefliste} sammelt diese Restbefunde nach der Bereinigung und der Zuordnung: 161 Befunde an 139 der 499 Kundensätze (\cref{tab:pruefliste}).

\begin{table}[H]
  \centering\small
  \caption{Restbefunde in \pq{CHECK\_Pruefliste} nach Bereinigung und Zuordnung}
  \label{tab:pruefliste}
  \begin{tabular}{@{}>{\raggedright\arraybackslash}p{6.4cm}r>{\raggedright\arraybackslash}p{7.0cm}@{}}
    \toprule
    Befund & Zeilen & Was manuell zu prüfen und zu entscheiden ist \\
    \midrule
    Adresse bei mehreren Kunden & 50 & Filiale oder Vereinsheim (getrennt lassen) oder derselbe Kunde nach Umzug (zusammenführen) \\
    Telefon enthält Text & 28 & Durchwahl und Ansprechpartner aus dem Resttext übernehmen \\
    E-Mail bei mehreren Kunden & 22 & Sammelpostfach bestätigen \\
    Erstkontakt nur monatsgenau & 12 & keine Korrektur möglich, Unschärfe ist markiert \\
    Name evtl.\ abgeschnitten (Feldlänge 30) & 9 & vollen Namen aus CRM oder Shop übernehmen \\
    PLZ fehlt, ungültig oder passt nicht zum Land & 7 & nachschlagen \\
    E-Mail ungültig oder Text im E-Mail-Feld & 6 & Adresse erfragen, Hinweis in die Bemerkung \\
    Privatname im Firmenfeld & 6 & Vor- und Nachname in die richtigen Felder \\
    Sperrvermerk im Freitext oder Firmenfeld & 5 & Sperre im dafür vorgesehenen Feld setzen \\
    Testkunde, zwei E-Mails, Postfach, Sonstige & 16 & zur Kenntnis; Testkunden sind bereits gefiltert \\
    \bottomrule
  \end{tabular}
\end{table}

Ein Befund verdient eine eigene Betrachtung. Hinter den 50 Zeilen \enquote{Adresse bei mehreren Kunden} stehen zwölf Adressen. Acht davon sind Vereinsheime: Verein und Gaststätte teilen sich das Haus und bleiben zu Recht getrennt. Die übrigen vier sind Umzieher. Bei ihnen trägt der Shop-Datensatz noch die alte Adresse und den Namen des ERP-Kunden, hängt aber über die E-Mail am neuen CRM-Eintrag -- der Shop verbindet also zwei Schlüssel, die die paarweise Kaskade getrennt gelassen hat. Eine Brückenregel könnte diese vier Fälle automatisch zusammenführen; im Nachbau der Logik sinkt die Zahl der zerrissenen Kunden damit von sieben auf drei, ohne neue Verschmelzung. Sie wurde bewusst \emph{nicht} eingebaut. Automatisch zusammenführen heißt auch automatisch entscheiden, welche Adresse, welcher Ansprechpartner und welcher Betreuer gelten, und eine falsche Zusammenführung überschreibt Daten, die sich aus dem integrierten Bestand nicht mehr zurückgewinnen lassen. Aus der Formseite sind ein Umzug und ein gemeinsam genutztes Gebäude nicht zu unterscheiden. Der richtige Ort für solche Fälle ist deshalb eine Abfrage, die die betroffenen Datensätze nebeneinander zeigt; sie müssen manuell überprüft und entschieden werden, erst danach wird zusammengeführt oder getrennt. Genau das leistet die Prüfliste.

\todo{Drei Absätze. (1) Kriterium: Bei 499 Kundensätzen lohnt eine Regel ab etwa 5 gleichartigen Fällen, darunter ist die Einzelprüfung schneller. Die Klasse hängt an drei Fragen: sicher erkennbar, eindeutig korrigierbar, Kosten einer falschen Korrektur. (2) Prüfliste erst nach der Bereinigung: sie enthält nur, was übrig bleibt. Zahlen aus \pq{CHECK\_Fehlerlog} einsetzen. (3) Was verworfen wird und was nicht: Testkunden, Dubletten, Stornos, Waisen, Doppelimporte -- gezählt; kein Datensatz wegen Qualitätsmängeln.}

\subsection{Fakten und Auswertung}\label{sec:fakten}

\todo{Hier begründen: Verkäufer aus dem Auftrag, Vertreter laut Kundenstamm als eigenes Attribut (F5); Shop-Bestellungen, die auch im ERP stehen, zählen nur einmal.}

\todo{Was: komposite Schlüssel verketten, Produktgruppe aus sprechendem Schlüssel und ArtNr$\leftrightarrow$SKU, CHF$\rightarrow$EUR, Brutto$\rightarrow$Netto je Land, Sets$\rightarrow$Stück, Filter (cancelled/refunded: 56 Bestellungen) mit Zählung, Auftrag $\rightarrow$ zum Datum gültige Adresse, Shop-Aufträge nur aus einer Quelle (ShopRef, 131 per Anti-Join). Fehler in den Bewegungsdaten: 4 Positionen ohne Auftrag (Inner Join, gezählt), 3 doppelt importierte Shop-Aufträge (je ShopRef der erste), 4 Lieferungen vor Auftragsdatum (Lieferdauer leer statt negativ), 5 Gutschriften (bleiben und mindern den Umsatz). UTC $\rightarrow$ MEZ vor dem Abschneiden der Uhrzeit.}


\abb{vn-7-bewegung}{Vorher und nachher, Gruppe 7: Bewegungsdaten und Kennzahlen. Verworfen werden nur Zeilen, die Umsatz doppelt zählen oder zu keinem Auftrag gehören}{fig:vn7}{1.0}
\abb{check-fakten}{\pq{CHECK\_Fakten}: Soll-Ist-Prüfung der Faktenseite -- Zeilenzahlen, Anti-Join, fehlende Schlüssel, Umsatzsummen je Quelle}{fig:checkfakten}{0.6}

\abb{dashboard}{Ergebnis-Dashboard: alle sechs Auswertungsfragen F1--F6 als PivotCharts auf dem integrierten Datenbestand (Blatt \pq{Dashboard}); jede Diagrammüberschrift formuliert die Kernaussage \skript{4.3}}{fig:dashboard}{1.0}

\todo{Dashboard kurz einordnen: eine Bildschirmseite, Diagramm statt Tabelle, Aussage als Titel, Zeitreihe als Linie, Struktur als Säule/Balken/Kreis (Skript Kap.~4.3). Chart-Wahl je Frage und was ohne Integration nicht darstellbar wäre (F5: Stamm-Vertreter aus CRM/ERP vs.\ Kanal aus Shop).}

\todo{Je Pivot zwei Sätze: was sie zeigt und welche Integrationsleistung sie voraussetzt.}

% =====================================================================
\section{Zusammenfassung}\label{sec:zusammenfassung}

\todo{Eine halbe Seite: drei Quellsysteme mit absichtlich eingebauten Integrationsproblemen, per Skript mit Sollwerten erzeugt; Power-Query-Mappe mit Staging, Harmonisierung, vierstufiger Kaskade, Sternschema, Prüfabfragen. Ergebnis: 285 Kunden aus 499 Kundensätzen mit 318 eingebauten Fehlern, 0 falsche Verschmelzungen, 7 bewusst getrennte Umzieher, Prüfliste mit Restbefunden; alle sechs Fragen per Pivot beantwortbar. Zentrale Entscheidungen: Historisierung statt Priorisierung, Verkäufer am Auftrag, Precision vor Recall, Fuzzy nur als Kandidatenlieferant mit fachlicher Nachregel. Kein Merkmal zählt allein. Eigenleistung: über das Skript hinausgehende Fehlerbilder, zwei Sollwert-Dateien, Messreihe als Methode.}

% =====================================================================
% ANHANG – knapp (HR 2.9: nie größer als der Textteil)
% =====================================================================
\anhang{Zentrale Power-Query-Abfragen (M-Code)}{app:mcode}
Der vollständige M-Code liegt in der Mappe; hier die Abfragen, die die Entscheidungen aus \cref{sec:integration} tragen.
% Zeilennummern passen zu code/Section1.m (Stand 12.11.2026); bei Änderungen in der Mappe neu extrahieren.
\lstinputlisting[language=PowerQueryM, caption={\pq{fnNormName}: Namen für den Abgleich normalisieren}, firstline=70, lastline=80]{code/Section1.m}
\lstinputlisting[language=PowerQueryM, caption={\pq{fnNameOk}: Namensbestätigung über den gemeinsamen Anfang}, firstline=108, lastline=116]{code/Section1.m}
\lstinputlisting[language=PowerQueryM, caption={\pq{CRM\_k}: Match-Kaskade}, firstline=355, lastline=404]{code/Section1.m}
\lstinputlisting[language=PowerQueryM, caption={\pq{CHECK\_Sollwerte}: Erfolgsmessung}, firstline=481, lastline=505]{code/Section1.m}

\anhang{Generator der Beispieldaten}{app:generator}
\texttt{generate\_sources.py} erzeugt alle Quelldateien deterministisch (Seed 20261127) und schreibt die Sollwerte, die nur der Erfolgsmessung dienen.
Das vollständige Skript (rund 540 Zeilen) liegt im Abgabe-Ordner; hier die Stellen, die den Entwurf tragen: Seed und skalierbare Fehlerquote, die Kundenverteilung auf die Systeme, die systemspezifischen Schreibweisen und das Schreiben der Sollwerte.
\lstinputlisting[language=Python, basicstyle=\ttfamily\scriptsize, caption={generate\_sources.py: Seed und Fehlerquote}, firstline=19, lastline=23]{code/generate_sources.py}
\lstinputlisting[language=Python, basicstyle=\ttfamily\scriptsize, caption={generate\_sources.py: Kunden und Verteilung auf die Quellsysteme}, firstline=83, lastline=113]{code/generate_sources.py}
\lstinputlisting[language=Python, basicstyle=\ttfamily\scriptsize, caption={generate\_sources.py: Schreibweisen je System}, firstline=144, lastline=170]{code/generate_sources.py}
\lstinputlisting[language=Python, basicstyle=\ttfamily\scriptsize, caption={generate\_sources.py: Sollwerte (Ground Truth)}, firstline=519, lastline=527]{code/generate_sources.py}

% ---------- Ehrenwörtliche Erklärung (HR 2.12, Wortlaut vorgegeben, ohne Seitenzahl)
\newpage
\thispagestyle{empty}
\anhang{Vorher-Nachher-Belege}{app:belege}

\abb{pq-beleg}{Die Abfrage \pq{BELEG\_Fehlerbilder} im Power-Query-Editor: 70 Beispiele in sieben Gruppen; links die Abfragen der Mappe}{fig:pqbeleg}{1.0}

\abb{vn-kunde-34}{Kunde 34 (Händler, Österreich) in den drei Quellen. Im CRM fehlt das Land; eine vierstellige PLZ lässt Österreich und die Schweiz zu, deshalb bleibt es dort leer. Der Ortszusatz des CRM bleibt stehen}{fig:vnkunde34}{1.0}

\section*{Ehrenwörtliche Erklärung}
Ich versichere hiermit, dass ich die vorliegende Arbeit mit dem Thema \enquote{\arbeitstitel} selbständig verfasst und keine anderen als die angegebenen Quellen und Hilfsmittel benutzt habe. Ich versichere zudem, dass die eingereichte elektronische Fassung mit der gedruckten Fassung übereinstimmt.

\vspace{2cm}
\makebox[0.5\textwidth]{\hrulefill} \hfill \makebox[0.4\textwidth]{\hrulefill} \\
Ort, Datum \hfill Unterschrift

\end{document}
